\section{Проверка}
\label{sec:verification}
\porozhdeno{\Claim{V01}{Проверка исходного сравнения отделена от его расчёта: она сверяет частоты и назначения с артефактами QIIME, полноту идентификаторов, роли образцов, суммы категорий и расстояния. Испытания используют малый пример с ручными ответами и искусственно повреждённые выходы; локальное воспроизведение сравнивается с сохранёнными таблицами. Для семейственных профилей независимо проверяются маски отбора, суммы и отображаемые группы, включая отказ при подмене рида.}}

\porozhdeno{\Claim{V02}{В трёх технических исполнениях DADA2 для всех \ProjectSamplesAll{} образцов, включая Kr5\_S84, попарно совпали наборы ASV-последовательностей, все частоты ASV по образцам и счётчики \texttt{input}, \texttt{filtered}, \texttt{denoised}, \texttt{non-chimeric}; это подтверждает повторяемость вычислений. Контроль качества обработки оценивается отдельно: сохранность Kr5\_S84 равна \ProjectKrFiveRetained{}/\ProjectKrFiveInput{} (\ProjectKrFiveRetentionPct{}\%), ниже порога \RevisionRetentionThresholdPct{}\%, поэтому полный верификатор возвращает \texttt{INCOMPLETE} с кодом завершения 2. Отдельная проверка повторяемости даёт \texttt{PASS}, допуская только эту согласованную потерю ридов; Kr5 остаётся в сравнении запусков и во всех опубликованных сводках.}}

\porozhdeno{\Claim{V03}{Перенос результатов в числовой комплект проверен независимо от генератора: выполнены \ArticleTransferChecks{} сверок источников, знаменателей, округления и значений TeX. В \ArticleTransferCases{} испытаниях командной строки обнаружены все \ArticleDetectedMutations{} содержательных подмен; исходный и восстановленный комплекты приняты. Эти проверки устанавливают соответствие сохранённым данным, а не истинную таксономию; смысл утверждений статьи дополнительно сопоставлен с первичными источниками и картой оснований.}}
